\chapter{The Construction of a Special Class of Automorphic Forms}

\label{ch:7}
\label{ch:7}

The problem of constructing automorphic forms is a difficult one. Of course one knows that holomorphic cusp forms of integral weight exist and one knows how to construct them using Poincar\'e series (cf.\ Example B(ii)). On the other hand one knows essentially nothing about which real analytic cusp forms exist let alone how to construct them. In this direction the pioneering work is due to Maass. Maass reduced the problem of constructing real analytic wave forms to the study of certain L-series with functional equation attached to real (as opposed to complex) quadratic fields. In his work (as well as in related works by Hecke) theta series attached to grossencharacters played a fundamental role.

The purpose of this section is to explain how Jacquet-Langland's theory considerably generalizes and smooths out the techniques of (Hecke and) Maass. The result is the construction of a wide (but still special) class of cusp-forms indexed by certain grossencharacters of global fields. The representations in question are constructed using the representation-theoretic formulation of theta-functions due to Weil. The fact that they actually define cusp forms will follow from a study of their L-functions.

In the first subsection below we summarize some salient features of the so-called Weil representation and its applicability to representation theory. We especially have in mind the problem of constructing supercuspidal representations. Because the Weil representation is also used in Section~\ref{ch:10} in a somewhat different context we develop it in greater generality than is required at present.

In Subsection B we apply the Weil representation to the construction of cusp forms as alluded to above. Then in Subsection C this construction is specialized to provide explicit examples of non-zero real analytic cusp forms of specified level $N$ and eigenvalue $s$. Subsection D sketches the connections between this construction (and its generalizations) and certain natural conjectures in non-abelian class field theory.

\section{The Weil Representation}\label{sec:7a}

Let $F$ denote a local field and $G$ the group $\GL(2,F)$. In this context the Weil representation is a representation of $G$ which one canonically associates to any quadratic form defined over $F$.

For our purposes it will suffice to consider forms $q$ of the following type:

(i) $q$ is the norm form of a separable quadratic extension $L$ of $F$; or

(ii) $q$ is the reduced norm of the unique quaternion division algebra $D$ over $F$.

In either case let $x \to x^{\sigma}$ denote the canonical involution of $V = L$ or $D$ over $F$. Then

\[
q(x) = xx^{\sigma}
\]

for all $x$ in $V$ and

\[
x + x^{\sigma} = \tr(x)
\]

where $(tr)$ denotes the trace (or reduced trace) from $V = L$ (or $D$) to $F$.

Now fix a non-trivial additive character $\tau$ of $F$. Since $(x,y) \to \tr(xy)$ is a non-degenerate bilinear form on $V$, $V$ can be identified with its dual by the pairing

\[
\langle x, y \rangle = \tau(\tr(xy)).
\]

The Fourier transform $\hat{\Phi}$ of each $\Phi$ in the Schwartz-Bruhat space $S(V)$ is by definition

\begin{equation}\label{eq:7.1}
\hat{\Phi}(x) = \int_{D} \Phi(y) \langle x, y \rangle \dd y\tag{7.1}
\end{equation}

where Haar measure $\dd y$ is normalized so that $(\hat{\Phi})^{\wedge}(x) = \Phi(-x)$.

To describe the Weil representation corresponding to the pair $(q,V)$ it is necessary to introduce an important invariant for $V$. First put

\begin{equation}\label{eq:7.2}
f(x) = \tau(q(x)).\tag{7.2}
\end{equation}

By \cite{Weil3} there exists a constant $\gamma$ depending only on $D$ (and $\tau$) such that

\begin{equation}\label{eq:7.3}
(\Phi^{*}f)^{\hat{}}(x) = \gamma f^{-1}(x^{\sigma})\hat{\Phi}(x)\tag{7.3}
\end{equation}

for all $\Phi \in S(V)$. In particular, $\gamma = -1$ if $V = D$. If $V = L$ suffice it to say that $|\gamma| = 1$.

Now consider the representation $r(s)$ of $\SL(2,F)$ in $S(V)$ such that

\begin{equation}\label{eq:7.4}
r\left(\begin{pmatrix} 1 & u \\ 0 & 1 \end{pmatrix}\right)\Phi(x) = \tau(uq(x))\Phi(x),\tag{7.4}
\end{equation}

\begin{equation}\label{eq:7.5}
r\left(\begin{pmatrix} a & 0 \\ 0 & a^{-1} \end{pmatrix}\right)\Phi(x) = \omega(a)|a|_V^{1/2}\Phi(ax),\tag{7.5}
\end{equation}

and

\begin{equation}\label{eq:7.6}
r\left(\begin{pmatrix} 0 & 1 \\ -1 & 0 \end{pmatrix}\right)\Phi(x) = \gamma\hat{\Phi}(x^{\sigma}).\tag{7.6}
\end{equation}

Here $\omega$ denotes the non-trivial character of $F^{\times}/\mathrm{Norm}(V^{\times})$ if $V = L$ and the trivial character of $F^{\times}$ if $V = D$. Since elements of the form

\begin{equation}\label{eq:7.7}
\begin{pmatrix} a & 0 \\ 0 & a^{-1} \end{pmatrix}, \quad \begin{pmatrix} 1 & u \\ 0 & 1 \end{pmatrix}, \quad \text{and} \quad \begin{pmatrix} 0 & 1 \\ -1 & 0 \end{pmatrix}\tag{7.7}
\end{equation}

generate $\SL(2,F)$ there will be at most one representation of $\SL(2,F)$ satisfying \eqref{eq:7.4}-\eqref{eq:7.6}. That there is one such representation (call it $r(s)$) is a result of Weil, Shalika, and Tanaka. To prove it one has simply to verify that the mapping

\[
s \to r(s)
\]

specified by \eqref{eq:7.4}-\eqref{eq:7.6} preserves all relations between the generators \eqref{eq:7.7}.

The representation $r(s)$ just described is ``the'' Weil representation of $\SL(2,F)$ associated to the pair $(V,q)$ and it may depend on the choice of character $\tau$. However, if $a \in F^{\times}$, and $\tau_a(x) = \tau(ax)$, then the representation $r_a(s)$ corresponding to $\tau_a$ is related to $r(s)$ through the formula

\begin{equation}\label{eq:7.8}
r_a(s) = r\left(\begin{pmatrix} a & 0 \\ 0 & 1 \end{pmatrix} s \begin{pmatrix} a^{-1} & 0 \\ 0 & 1 \end{pmatrix}\right).\tag{7.8}
\end{equation}

If $b \in V^{\times}$ define operators $\lambda(b)$ and $\rho(b)$ in $S(V)$ by

\begin{equation}\label{eq:7.9}
\lambda(b)\Phi(x) = \Phi(b^{-1}x)\tag{7.9}
\end{equation}

and

\begin{equation}\label{eq:7.10}
\rho(b)\Phi(x) = \Phi(xb).\tag{7.10}
\end{equation}

Then if $a = q(b)$, a straightforward computation shows that

\begin{equation}\label{eq:7.11}
r_a(s)\lambda(b^{-1}) = \lambda(b^{-1})r(s)\tag{7.11}
\end{equation}

and

\begin{equation}\label{eq:7.12}
r_a(s)\rho(b) = \rho(b)r(s).\tag{7.12}
\end{equation}

In particular $r(s)$ and $r_a(s)$ are equivalent if $a$ belongs to the subgroup $q(V^{\times})$ of $F^{\times}$. (So since $q(V^{\times}) = F^{\times}$ if $V = D$ and $F$ is not $\R$, the Weil representation in this case is independent of the choice of $\tau$.) Furthermore from \eqref{eq:7.11} and \eqref{eq:7.12} it is clear that both $\lambda(b)$ and $\rho(b)$ commute with $r(s)$ if $q(b) = 1$.

\begin{remark}\label{rem:7.1}

If $S(V)$ is given its usual topology, $r(s)$ is continuous. Furthermore $r(s)$ extends to a unitary representation of $\SL(2,F)$ on $\bL^2(V)$.

\end{remark}

\begin{remark}\label{rem:7.2}

(Concerning the General Construction of Weil.) The representations $r(s)$ corresponding to the forms $q$ on $L$ and $D$ represent special cases of a very general construction devised by Weil, Segal and Mackey. In Weil's theory a ``big'' representation is canonically attached to the ``abstract symplectic group'' of any locally compact abelian group $G$. The connection between this set-up and the above representations is that for certain special choices of $G$ the group $\SL(2,F)$ imbeds naturally in the symplectic group of $G$ and $r(s)$ appears as the restriction of this big representation to $\SL(2,F)$.

In rough detail, the situation is as follows. Fix a duality $\langle \cdot, \cdot \rangle$ between $G$ and its dual $G^*$. For each $w = (u, u^*)$ in $G \times G^*$, let $U(w)$ be the unitary operator in $\bL^2(G)$ taking $\phi(x)$ to $\phi'(x) = \langle x, u^*\rangle\phi(x + u)$. Then $U(w)$ is a projective representation of $G \times G^*$ with multiplier

\[
F(w_1, w_2) = \langle u_1, u_2^*\rangle \text{ if } w_i = (u_i, u_i^*), \text{ i.e.\ } U(w_1)U(w_2) = U(w_1 + w_2)F(w_1, w_2).
\]

Let $\overline{A}(G)$ denote the group of all unitary operators in $\bL^2(G)$ of the form $tU(w)$ with $|t| = 1$ and $\overline{B}_0(G)$ its normalizer in $U(\bL^2(G))$. Since $\overline{A}(G)$ is isomorphic to the group $A(G)$ of all pairs $w, t$ with $w \in G \times G^*$, $|t| = 1$, and $(w_1, t_1)(w_2, t_2) = (w_1 + w_2, F(w_1, w_2)t_1t_2)$, each member of $\overline{B}_0(G)$ induces an automorphism of $A(G)$ which has the form $(w, t) \to (w\sigma, tg(w))$ where $\sigma$ is an automorphism of $G \times G^*$ and $g$ is a continuous function from $G \times G^*$ to the circle group related to $\sigma$ by the identity

\[
F(w_1\sigma, w_2\sigma) = F(w_1, w_2)g(w_1 + w_2)/g(w_1)g(w_2).
\]

Conversely any pair $\sigma, g$ with $\sigma$ and $g$ so related defines an automorphism of $A(G)$ and the group of all such is denoted by $B_0(G)$.

Now let $p_0$ denote the homomorphism of $\overline{B}_0(G)$ into $B_0(G)$ just described. Weil proves that $p_0$ is surjective with kernel equal to the group of constant multiples of the identity. Consequently one obtains a projective representation of $B_0(G)$ in $\bL^2(G)$. The symplectic group $\operatorname{Sp}(G)$ associated to $G$ is the group of automorphisms of $G \times G^*$ preserving the form $F(w_1, w_2)F(w_2, w_1)^{-1}$. Equivalently $\sigma = \begin{pmatrix} a & \beta \\ \gamma & \delta \end{pmatrix}$ is symplectic if $\sigma\sigma^I = I$ where $\sigma^I = \begin{pmatrix} \delta^* & -\beta^* \\ -\gamma^* & a^* \end{pmatrix}$. (Here $a, \beta, \gamma$ and $\delta$ are homomorphisms of $G$ into $G$, of $G$ into $G^*$, of $G^*$ into $G$, and $G$ into $G^*$ respectively, $w\sigma = (u, u^*)\sigma = (u, u^*)\begin{pmatrix} a & \beta \\ \gamma & \delta \end{pmatrix}$, and $a^*$ denotes the dual of $a$.) Since $\operatorname{Sp}(G)$ injects in $B_0(G)$ through the homomorphism $\begin{pmatrix} a & \beta \\ \gamma & \delta \end{pmatrix} \to (\sigma, g_\sigma)$ where

\[
g_\sigma(u, u^*) = \langle u, 2^{-1}u\alpha\beta^*\rangle \langle 2^{-1}u^*\gamma\delta^*, u^*\rangle \langle u^*\gamma, u\beta\rangle
\]

there is associated to $\operatorname{Sp}(G)$ a natural (possibly projective) representation in $\bL^2(G)$. This is the Weil representation attached to $G$.

To obtain the representation $r(s)$ of Remark~\ref{rem:7.1} take $G$ to be the additive group of $V$ and fix the self-duality

\[
\langle u, v \rangle = \tau(2B_q(u, v))
\]

if $B_q$ denotes the bilinear form associated to $(q,V)$. Since each $a \in F$ defines a homomorphism of $G = V^+$ (by scalar multiplication), and since $\begin{pmatrix} a & \beta \\ \gamma & \delta \end{pmatrix}$ in $\SL(2,F)$ satisfies $\sigma\sigma^I = I$, each element of $\SL(2,F)$ imbeds homomorphically in $B_0(G)$. By restriction there results an (ordinary) representation of $\SL(2,F)$ which coincides with the special representation $r(s)$ described earlier.

Before returning (finally) to our discussion for $\SL(2,F)$ let us note that the assertion ``$p_0$ is surjective'' is essentially an analogue (for local fields) of the Stone-von Neumann Uniqueness Theorem for representations of the Heisenberg commutation relations. Indeed this theorem asserts that (up to equivalence) there is exactly one irreducible representation of $G \times G^*$ with multiplier $F$.

\end{remark}

To describe the Weil representation for $\GL(2,F)$ corresponding to the pair $(q,V)$ one introduces the group $G_+$ consisting of $n$ in $\GL(2,F)$ with determinant in $q(V^{\times})$. This group has index 2 or 1 in $\GL(2,F)$ according if $V = L$ or $V = D$ (with $F \neq \R$).

\begin{lemma}\label{lem:7.3}

The representation $r(s)$ of $\SL(2,F)$ on $S(V)$ extends to a representation $r(n)$ of $G_+$ such that

\begin{equation}\label{eq:7.13}
r\left(\begin{pmatrix} a & 0 \\ 0 & 1 \end{pmatrix}\right)\Phi(x) = \Phi(xh)\tag{7.13}
\end{equation}

if $a = q(h)$ belongs to $q(V^{\times})$.

\end{lemma}

\begin{proof}

Let $C$ denote the group of matrices

\[
\left\{\begin{pmatrix} a & 0 \\ 0 & 1 \end{pmatrix}\right\}
\]

with $a$ in $q(V^{\times})$ and observe that $G_+$ is the semi-direct product of $G$ and $\SL(2,F)$. Clearly \eqref{eq:7.13} defines a continuous representation of $H$ on $S(V^{\times})$. Therefore to prove the lemma it suffices to verify the relation

\begin{equation}\label{eq:7.14}
r\left(\begin{pmatrix} a & 0 \\ 0 & 1 \end{pmatrix} s \begin{pmatrix} a^{-1} & 0 \\ 0 & 1 \end{pmatrix}\right) = r\left(\begin{pmatrix} a & 0 \\ 0 & 1 \end{pmatrix}\right)r(s)r\left(\begin{pmatrix} a^{-1} & 0 \\ 0 & 1 \end{pmatrix}\right).\tag{7.14}
\end{equation}

But the left side of \eqref{eq:7.14} is simply $r_a(s)$ and

\[
r\left(\begin{pmatrix} a & 0 \\ 0 & 1 \end{pmatrix}\right)\Phi(x) = \rho(h)\Phi(x).
\]

Consequently \eqref{eq:7.14} and the lemma follow from \eqref{eq:7.12}.

\end{proof}

To emphasize the possible dependence of $r(n)$ on $\tau$ we denote it by $r_\tau(n)$. Now let $r(g)$ denote the representation of $\GL(2,F)$ induced from $r_\tau(n)$. This $r(g)$ is ``the Weil representation of $\GL(2,F)$ (associated to the pair $(q,V)$) and it is independent of the choice of character $\tau$. (Indeed $r(g)$ is obtained by letting $\GL(2,F)$ act by right translation on the space of functions $\phi$ on $G$ (with values in the space of $r_\tau(n)$) which satisfy

\[
\phi(ng) = r_\tau(n)\phi(g)
\]

for all $n \in G_+$. If $a \in F^{\times}$ and we define $\phi^{a}(g) = \phi\left(\begin{pmatrix} a & 0 \\ 0 & 1 \end{pmatrix}g\right)$ then $\phi^{a}$ is 0 iff $\phi$ is, and

\[
\begin{aligned}
\phi^{a}(ng) &= \phi\left(\begin{pmatrix} a & 0 \\ 0 & 1 \end{pmatrix}ng\right)\\
&= r_\tau\left(\begin{pmatrix} a & 0 \\ 0 & 1 \end{pmatrix}n\begin{pmatrix} a^{-1} & 0 \\ 0 & 1 \end{pmatrix}\right)\phi^{a}(g).
\end{aligned}
\]

That is, replacing the space $\{\phi\}$ by the space $\{\phi^a\}$, we obtain an equivalent representation of $G$ induced from the representation

\[
n \to (r_\tau)_a(n)
\]

of $G_+$. So since $a$ is arbitrary it follows that $r(g)$ is independent of $\tau$.)

We close this subsection by describing how the representations already introduced allow one to construct a wide class of irreducible admissible representations of $\GL(2,F)$. This connection between Weil's construction and irreducible representation theory was first made explicit in \cite{Shalika} and \cite{Tanaka} and then completely exploited in \cite{JacquetLanglands}. The crucial point is a relation of the type \eqref{eq:7.11} or \eqref{eq:7.12}. Such a relation asserts that the regular representation of $V^{\times}$ in $S(V)$ commutes with $r(s)$. This in turn implies that the Weil representation is highly reducible and in fact decomposes into irreducibles according to the representation theory of $V^{\times}$.

So suppose $\pi'$ is an irreducible (finite-dimensional) representation of $V^{\times} = G'$ on a vector space $H$ defined over $\C$. Taking the tensor product of $r(s)$ with the trivial representation of $\SL(2,F)$ on $H$ one obtains a representation on

\[
S(V) \otimes_{\C} H
\]

which we again denote by $r$. An element in $S(V) \otimes H$ will be regarded as a function on $V$ taking values in $H$ whose coordinate entries (with respect to some basis of $H$) are Schwartz-Bruhat functions on $V$. Because of \eqref{eq:7.12} the subspace of functions in $S(V) \otimes H$ satisfying

\begin{equation}\label{eq:7.15}
\Phi(xh) = \pi'(h^{-1})\Phi(x)\tag{7.15}
\end{equation}

for all $h$ in $G'$ with $q(h) = 1$ is immediately seen to be invariant for $r(s)$. The resulting subrepresentation of $r(s)$ will be denoted by $r_{\pi'}(s)$.

By the reasoning used to prove Lemma~\ref{lem:7.3} it can be shown that $r_{\pi'}(s)$ extends to a representation of $G_+$ satisfying

\begin{equation}\label{eq:7.16}
r_{\pi'}\left(\begin{pmatrix} a & 0 \\ 0 & 1 \end{pmatrix}\right)\Phi(x) = |h|_V^{1/2}\pi'(h)\Phi(xh)\tag{7.16}
\end{equation}

if $a = q(h)$ belongs to $q(V^{\times})$. Furthermore

\[
r_{\pi'}\left(\begin{pmatrix} a & 0 \\ 0 & a \end{pmatrix}\right) = \omega(a)\chi_{\pi'}(a)I
\]

if $a$ is arbitrary in $F^{\times}$ and $\chi_{\pi'}$ is the central character of $\pi'$. If $\pi'$ is unitary (and $H$ is a Hilbert space) then $r_{\pi'}$ can be extended to a unitary representation of $G_+$ in $\bL^2(V,\pi')$, the closure of $S(V,\pi')$ in the Hilbert space of square integrable functions from $V$ to $H$ with norm $\|\Phi\|^2 = \int_V \|\Phi(x)\|^2 \dd x$.

Now consider the representation of $\GL(2,F)$ induced from $r_{\pi'}(n)$ and denote it again by $r_{\pi'}$. Of course $r_{\pi'}(g)$ is independent of the choice of character $\tau$ since it is a component of the Weil representation $r(g)$ of $\GL(2,F)$. What is remarkable is that $r_{\pi'}(g)$ is always admissible and (essentially) always an irreducible representation. The equivalence class of $r_{\pi'}(g)$ depends on the choice of $\pi'$ and on whether $V$ equals $L$ or $D$.

\begin{theorem}\label{thm:7.4}

($V$ is a separable quadratic extension of $F$.)

\begin{enumerate}[(i)]

\item If there is no character $\delta$ of $F^{\times}$ such that $\pi' = \delta \circ q$ then $r_{\pi'}(g)$ is a super cuspidal representation of $\GL(2,F)$.

In fact by varying such $\pi'$ as $V$ varies through all equivalence classes of quadratic extensions, all super cuspidal representations of $\GL(2,F)$ may be so obtained;

\item If $\pi' = \delta \circ q$ for some character $\delta$ of $F^{\times}$ then $r_{\pi'}(g)$ is equivalent to the principal series representation $\pi(\delta, \delta\omega)$.

\end{enumerate}

\end{theorem}

\begin{remark}\label{rem:7.5}

The kernel of $q$ is the norm one group $L^1$ of $L$. Thus $\pi' = \delta \circ q$ for some character $\delta$ of $F^{\times}$ if and only if $\pi'$ is trivial on $L^1$. In particular, $r_{\pi'}$ is super cuspidal only if $\pi'$ is a ramified character of $L^{\times}$. Note too that if $\pi'$ is trivial on $L^1$ it factors through $F^{\times}$ by exactly two characters, namely $\delta$ and $\delta\omega$.

\end{remark}

\begin{theorem}\label{thm:7.6}

($V$ is the unique division quaternion algebra $D$ defined over $F$.)

\begin{enumerate}[(i)]

\item The representation $r_{\pi'}(g)$ decomposes as the direct sum of $d = \dim(\pi')$ mutually equivalent irreducible representations $\pi(\pi')$ of $\GL(2,F)$;

\item Each $\pi(\pi')$ is absolutely cuspidal if $d > 1$ and special if $d = 1$;

\item All the super cuspidal and special representations of $\GL(2,F)$ are obtained in this manner. More precisely, the map

\[
\pi' \to \pi(\pi')
\]

gives a one-to-one correspondence between the equivalence classes of (finite-dimensional) irreducible representations of $G'$ and the equivalence classes of special and super cuspidal representations of $\GL(2,F)$.

\end{enumerate}

\end{theorem}

\begin{remark}\label{rem:7.7}

In stating the theorems above we have implicitly assumed $F$ to be non-archimedean. In case $F = \R$, let $D$ denote the quaternion algebra over $\R$ identified with the set of matrices of the form $\begin{pmatrix} a & b \\ -b & \overline{a} \end{pmatrix}$ with $a, b$ in $\C$. Then $q(h) = \det(h)$, identifying

\[
h = a_0 + a_1i + a_2j + a_3k
\]

with

\begin{equation}\label{eq:7.17}
h = \begin{pmatrix} a_0 + a_1\sqrt{-1} & a_2 + a_3\sqrt{-1} \\ -(a_2 - a_3\sqrt{-1}) & a_0 - a_1\sqrt{-1} \end{pmatrix}.\tag{7.17}
\end{equation}

Every irreducible representation of $G' = D^{\times}$ is known to be of the form

\begin{equation}\label{eq:7.18}
\pi'(h) = q(h)^{r}\rho_n(h)\tag{7.18}
\end{equation}

where $r \in \C$, and $\rho_n(h)$ is the $n$-th symmetric tensor power of the standard two dimensional representation $\rho_1$ of $G'$ associated with the identification \eqref{eq:7.17}. Then

\begin{equation}\label{eq:7.19}
\pi(\pi') = \sigma(\mu_1, \mu_2),\tag{7.19}
\end{equation}

where $\mu_1(a) = |a|^{r + n + \frac{1}{2}}$, and $\mu_2(a) = |a|^{r - \frac{1}{2}}[\sgn(a)]^{n}$, and $\sigma(\mu_1, \mu_2)$ denotes the appropriate discrete series representation of $\GL(2,\R)$ in the notation of Section~\ref{ch:4}. Furthermore every such representation of $\GL(2,\R)$ is obviously obtained.

In case $F = \R$ and $V = L = \C$ the statement of Theorem~\ref{thm:7.4} is simplest of all. If $\pi'(z)$ is a character of $\C^{\times}$ not of the form $\chi(z\overline{z})$ with $\chi$ a quasi-character of $\R^{\times}$, then $\pi(\pi')$ is of the form

\begin{equation}\label{eq:7.20}
\pi'(z) = (z\overline{z})^{r}z^{m}\overline{z}^{n}\tag{7.20}
\end{equation}

with $r \in \C$, and $m$ and $n$ two integers, one zero and the other positive. In this case, $r_{\pi'}$ is equivalent to the discrete series representation $\sigma(\mu_1, \mu_2)$, with $\mu_1\mu_2(t) = |t|^{2r}t^{m+n}\sgn(t)$ and $\mu_1\mu_2^{-1}(t) = t^{m+n}\sgn(t)$. If, on the other hand, $\pi'$ is of the form \eqref{eq:7.20} with $n + m = 0$, i.e.\ if $\pi'(z)$ is of the form $\chi(z\overline{z})$, then $r_{\pi'}$ is equivalent to the principal series representation $\pi(\mu_1, \mu_1\sgn(t))$.

\end{remark}

\textbf{SUMMING UP.} Let $V$ denote a division algebra over the local field $F$ and $q$ a quadratic form on $V$. To the pair $(q,V)$ there is associated a natural representation $r$ of $\GL(2,F)$ called the Weil representation of $G$. In case $V$ is a separable quadratic extension of $F$ and $q$ the norm form of this extension, the decomposition of $r(g)$ according to the irreducible representations $\pi'$ of $V^{\times}$ contains supercuspidal and principal series representations. In case $V$ is the unique division quaternion algebra over $F$ and $q$ its norm form, the decomposition of $r(g)$ according to the irreducible representations $\pi'$ of $V^{\times}$ contains all the supercuspidal and special representations of $\GL(2,F)$. (In case $F = \R$ replace ``supercuspidal'' and ``special'' by ``discrete series''.) In either case there results a correspondence

\begin{equation}\label{eq:7.21}
\pi' \to \pi(\pi')\tag{7.21}
\end{equation}

between irreducible representations of $G' = V^{\times}$ and $G = \GL(2,F)$.

In this section we shall henceforth restrict our attention to the correspondence \eqref{eq:7.21} for $G' = L^{\times}$. The case of quaternion algebras will play a fundamental role in Section~\ref{ch:10}. In this case the Weil representation corresponding to $(D,q)$ will play the role of (constructing) theta series associated to quaternary (as opposed to binary) quadratic forms.

\section{The Construction of Certain Special Representations of \texorpdfstring{$\GL(2,\A)$}{GL(2,A)}}\label{sec:7b}

Let $F$ denote a global field and $L$ a separable quadratic extension of $F$. By the main theorem of Jacquet-Langlands theory the problem of constructing (generalized) cusp forms on $\GL(2)$ is reduced to the problem of constructing representations of $\GL(2,\A)$ whose associated L-functions satisfy certain prescribed analytic conditions. The purpose of this subsection is to explain how the Weil representation allows one to construct a wide class of such representations indexed by grossencharacters of $L$.

If $\lambda$ is any grossencharacter of $L$ write

\begin{equation}\label{eq:7.22}
\lambda = \prod_w \lambda_w\tag{7.22}
\end{equation}

where each $\lambda_w$ is a character of $(L_w)^{\times}$ and $w$ is a place of $L$. We shall attach to $\lambda$ a representation

\[
\pi(\lambda) = \bigotimes_v \pi_v
\]

of $\GL(2,\A(F))$ by defining $\pi_v$ for each place $v$ of $F$ in terms of the character (or characters) $\lambda_w$ lying over $v$. This must of course be done in such a way that $\pi_v$ is class 1 for almost every $v$.

Suppose first that $v$ is a prime in $F$ which splits in $L$ with divisors $w$ and $w'$. Then $L_w$ and $L_{w'}$ are isomorphic to $F_v$ and $\lambda_w$ and $\lambda_{w'}$ may be viewed as characters of $F_v^{\times}$. In this case we set $\pi_v$ equal to the principal series (actually continuous series) representation $\pi(\lambda_w, \lambda_{w'})$.

Suppose, on the other hand, that $v$ does not split in $L$ but lies under the single prime $w$. Then $L_w$ is a genuine quadratic extension of $F_v$ and $\lambda_w$ is a character of $L_w^{\times}$. In this case the Weil construction

\[
\pi' \to \pi(\pi')
\]

is applicable with $\lambda_w$ playing the role of $\pi'$. Therefore we set $\pi_v$ equal to $\pi(\lambda_w)$.

\begin{proposition}\label{prop:7.8}

Let the family $\{\pi_v\}$ be as above. Then the formula

\begin{equation}\label{eq:7.23}
\pi(\lambda) = \bigotimes_v \pi_v\tag{7.23}
\end{equation}

defines an irreducible unitary representation of $\GL(2,\A(F))$.

\end{proposition}

\begin{proof}

It suffices to show that $\pi_v$ is class 1 for almost every $v$. But by Theorem~\ref{thm:4.23}, $\pi_v$ is not class 1 only if one of the following possibilities obtains:

(a) $\pi_v = \pi(\lambda_w)$ is absolutely cuspidal;

(b) $\pi_v = \pi(\lambda_w, \lambda_{w'})$ with $\lambda_w$ or $\lambda_{w'}$ ramified; or

(c) $\pi_v = \pi(\delta_v, \delta_v\omega_v)$ with $\delta_v$ or $\omega_v$ ramified.

Now by Remark~\ref{rem:7.5}, (a) obtains only if $\lambda_w$ is ramified, and (c) obtains only if $v$ ramifies in $L$ or if $\lambda_w = \delta_v \circ q_v$ is ramified. Since these latter possibilities occur for only finitely many places the proof is complete.

\end{proof}

To show that $\pi(\lambda)$ is a cusp form we must analyze its corresponding L-function. So let $L(s,\pi(\lambda))$ denote the L-function attached to $\pi(\lambda)$ by Jacquet-Langlands' theory. To prove that $L(s,\pi(\lambda))$ satisfies the conditions of Theorem~\ref{thm:6.18} (and hence that $\pi(\lambda)$ occurs in some $\bR_0^\psi$) we must show that

\begin{equation}\label{eq:7.24}
L(s, \chi \otimes \pi(\lambda) = \prod_v L(s, \chi_v \otimes \pi_v)\tag{7.24}
\end{equation}

is entire, bounded in vertical strips, and such that the functional equation \eqref{eq:6.44} obtains for every grossencharacter $\chi$ of $F$.

To this end we recall the idele norm map from the idele class group of $L$ to the idele class group of $F$. This map is defined by

\begin{equation}\label{eq:7.25}
N(\beta) = N((\beta_w)) = (a_v) = a.\tag{7.25}
\end{equation}

if $a_v = \prod_{w|v} N(\beta_w)$. Through it we can extend the grossencharacter $\chi$ of $F$ to a grossencharacter $\prod_w \chi_w$ of $L$.

Now let $L(s,\lambda\chi)$ denote the Hecke L-series attached to the ``automorphic form'' $\lambda\chi$ of $L^{\times} = \GL(1)$. The analytic properties of this L-function are well known. Therefore our plan of attack will be to compare $L(s,\chi \otimes \pi(\lambda))$ with $L(s,\lambda\chi)$. Recall that

\begin{equation}\label{eq:7.26}
L(s,\lambda\chi) = \prod_w L(s,\lambda_w\chi_w)\tag{7.26}
\end{equation}

where

\begin{equation}\label{eq:7.27}
L(s,\lambda_w\chi_w) = (1 - \lambda_w\chi_w(\tilde{\omega}_w)N(\tilde{\omega}_w)^{-s})^{-1}\tag{7.27}
\end{equation}

if $\lambda_w\chi_w$ is unramified, and equals one otherwise. Actually in \eqref{eq:7.26} the product is taken over the infinite places of $L$ as well. If $L_w = \R$,

\begin{equation}\label{eq:7.28}
L(s,\lambda_w\chi_w) = \pi^{-1/2(s + it_w + m_w)}\Gamma\left(\frac{s + it_w + m_w}{2}\right)\tag{7.28}
\end{equation}

for $\lambda_w\chi_w(x) = [\sgn(x)]^{m_w}|x|^{it_w}$, $m_w = 0$ or 1. If $L_w = \C$,

\begin{equation}\label{eq:7.29}
L(s,\lambda_w\chi_w) = (2\pi)^{-(s + t_w + m_w + n_w)}\Gamma(s + t_w + m_w + n_w)\tag{7.29}
\end{equation}

for $\lambda_w\chi_w(z) = |z|^{it_w}z^{m}\overline{z}^{n}$, $\inf(n,m) = 0$.

\begin{lemma}\label{lem:7.9}

\begin{equation}\label{eq:7.30}
L(s,\chi \otimes \pi(\lambda)) = L(s,\lambda\chi)\tag{7.30}
\end{equation}

\end{lemma}

\begin{proof}

As both sides of \eqref{eq:7.30} are defined by infinite products it will suffice to verify the equality ``place by place.'' Caution: The left side is a product over the primes of $L$ (with Euler factors of degree at most 1) and the right side is a product over the primes of $F$ (with Euler factors of degree at most 2).

Suppose first that $v$ splits in $L$ and lies under $w$ and $w'$. Then $\pi_v = \pi(\lambda_w, \lambda_{w'})$, and so by Theorem~\ref{thm:6.15},

\begin{equation}\label{eq:7.31}
L(s,\chi_v \otimes \pi_v) = L(s,\chi_w\lambda_w)L(s,\chi_{w'}\lambda_{w'}).\tag{7.31}
\end{equation}

if $v$ is finite. If $v$ is infinite the same result holds by Theorem~\ref{thm:6.16}.

Next suppose that $v$ is inert and $w|v$. Then $L_w$ is an unramified quadratic extension of $F_v$. If $\chi_w\lambda_w$ is ramified then $L(s,\chi_v \otimes \pi_v)$ is 1 whether $\pi_v$ is supercuspidal or not. Therefore the desired equality

\begin{equation}\label{eq:7.32}
L(s,\chi_v \otimes \pi_v) = L(s,\chi_w\lambda_w)\tag{7.32}
\end{equation}

obtains by default. If, on the other hand, $\chi_w\lambda_w$ is unramified, then $\pi_v$ cannot be supercuspidal by Remark~\ref{rem:7.5}. Since $N(\tilde{\omega}_w) = \tilde{\omega}_v^{2}$ (where $\tilde{\omega}_w$ denotes a local uniformizing variable at $w$)

\[
\begin{aligned}
L(s,\chi_w\lambda_w) &= (1 - \chi_w\lambda_w(\tilde{\omega}_w)N(\tilde{\omega}_w)^{-s})^{-1}\\
&= (1 - \delta_v(\tilde{\omega}_v^{2})\chi_v(\tilde{\omega}_v^{2})N(\tilde{\omega}_v)^{-2s})^{-1}\\
&= (1 - \delta_v(\tilde{\omega}_v)\chi_v(\tilde{\omega}_v)N(\tilde{\omega}_v)^{-s})^{-1}(1 + \delta_v(\tilde{\omega}_v)\chi_v(\tilde{\omega}_v)N(\tilde{\omega}_v)^{-s})^{-1}
\end{aligned}
\]

if $\lambda_w = \delta_v \circ q$. Since $v$ is inert, $\omega_v(\tilde{\omega}_v) = -1$.

Thus

\[
\begin{aligned}
&L(s,\chi_w\lambda_w) = (1 - \delta_v(\tilde{\omega}_v)\chi_v(\tilde{\omega}_v)N(\tilde{\omega}_v)^{-s})^{-1}(1 - \delta_v(\tilde{\omega}_v)\omega_v\chi_v(\tilde{\omega}_v)N(\tilde{\omega}_v)^{-s})^{-1}\\
&\qquad = L(s,\chi_v \otimes \pi(\lambda_v, \lambda_v\omega_v))\\
&\qquad = L(s,\chi_v \otimes \pi_v)
\end{aligned}
\]

and \eqref{eq:7.32} again obtains.

The case when $v$ is infinite or ramified is left to the reader.

\end{proof}

\begin{lemma}\label{lem:7.10}

$L(s,\chi \otimes \pi(\lambda))$ is entire, bounded in vertical strips, and satisfies the requisite functional equation for all $\chi$ provided there is no grossencharacter $\mu$ of $F$ such that $\lambda = \mu \circ N$ with $N$ as in \eqref{eq:7.25}.

\end{lemma}

\begin{proof}

By Lemma~\ref{lem:7.9},

\begin{equation}\label{eq:7.33}
L(s,\chi \otimes \pi(\lambda)) = L(s,\chi\lambda).\tag{7.33}
\end{equation}

Therefore $L(s,\chi \otimes \pi(\lambda))$ satisfies the requisite functional equation since $L(s,\chi\lambda)$ does by Hecke's abelian theory. (Thus far we have made no use of the hypothesis on $\lambda$.) It remains however to check that $L(s,\chi\lambda)$ is entire and bounded in vertical strips.

According to Hecke's theorem for $L(s,\chi\lambda)$ as defined by \eqref{eq:7.26} (cf.\ \cite{Weil4}, Theorem 5, p. 133) $L(s,\chi\lambda)$ is entire if $\omega\lambda$ is a ``non-trivial'' grossencharacter of $L$, i.e.\ if

\begin{equation}\label{eq:7.34}
\chi\lambda(a) \not\equiv |a|^{r} \quad \text{for} \quad a \in \A^{\times}(L).\tag{7.34}
\end{equation}

But \eqref{eq:7.34} is equivalent to the condition

\[
\lambda(a) \not\equiv \chi(a)|a|^{r}.
\]

So since $\chi$ and $|\cdot|^{r}$ both come from characters of the idele class group of $F$ ($\chi$ by definition and $|a|^{r} = \prod_w \psi_w$ because $\psi_w = \psi_{w'}$ whenever $w, w'|v$) it follows that condition \eqref{eq:7.34} is implied by our hypothesis on $\lambda$. Consequently $L(s,\chi\lambda)$ (and hence $L(s,\chi \otimes \pi(\lambda))$) is entire.

That $L(s,\chi\lambda)$ is bounded in vertical strips follows from the integral representation for the global zeta-function $\zeta(f,\chi\lambda,s)$ and the fact that $\zeta(f,\chi\lambda,s)$ essentially equals $L(s,\chi\lambda)$ for an appropriate choice of $f$.

\end{proof}

Summing up, we have:

\begin{theorem}\label{thm:7.11}

Let $L$ denote any separable quadratic extension of a global field $F$. The map

\[
\lambda \to \pi(\lambda)
\]

just described defines a correspondence between grossencharacters of $L$ (i.e.\ automorphic forms on $\GL(1)$ over $L$) and irreducible representations of the adele group of $\GL(2)$ over $F$. If $\lambda$ is not of the form $\mu \circ N$ for any grossencharacter $\mu$ of $F$ then $\pi(\lambda)$ is actually a cusp form on $\GL(2)$.

\end{theorem}

Suppose $\psi$ denotes the central character of $\pi(\lambda)$. From the theorem above it follows that there is a subspace of functions in $\bL_0^2(\bG_F \setminus \bG_A, \psi)$ which is equivalent as a $\bG_A$-module to $\pi(\lambda)$. In particular, if $F = \Q$, these functions describe classical holomorphic or real analytic cusp forms of level equal to the conductor of $\pi(\lambda)$. The drawback of the elegant proof of Theorem~\ref{thm:7.11} just sketched is that it does not construct these functions directly. This situation is in contrast to the classical approach involving theta series alluded to earlier. Since the constructive theme of this classical approach is desirable from several points of view we shall briefly sketch its modern formulation below.

The direct construction of cusp forms in the space of $\pi(\lambda)$ (in the generality of Theorem~\ref{thm:7.11}) is due to \cite{ShalikaTanaka}. Their point of departure is a global analogue of the Weil representation which may be described as follows. (Cf.\ \cite{Weil3}, Chapter III.)

As before let $L$ denote a separable quadratic extension of $F$, and $\tau = \prod_v \tau_v$ an additive character of $F$. If the place $v$ of $F$ lies under the single prime $w$ of $L$, let $L_v$ denote the completion $L_w$. On the other hand, if both $w$ and $w'$ in $L$ divide $v$, put $L_v$ equal to the direct sum $L_w \oplus L_{w'}$. In either case, $L_v$ will be a semi-simple commutative algebra of rank two over $F$ equipped with a quadratic form $q_v$ given by the relative norm from $L_v$ to $F_v$. (In case $L_v$ is a quadratic extension of $F_v$, $q_v$ is the usual norm form; in case $L_v = L_w \oplus L_{w'}$, $q_v = xx^{\sigma}$ with $x^{\sigma} = (b, a)$ if $x = (a, b)$.)

By the general theory of Weil there is a constant $\gamma_v$ (of absolute value 1) such that the formulas \eqref{eq:7.4}-\eqref{eq:7.6} define a representation of $\SL(2,F_v)$ on $S(L_v)$ (the local Weil representation). This representation extends to a unitary representation $r_v$ of $\bL^2(L_v)$ which coincides with the Weil representation introduced in Subsection A if $L_v$ is a quadratic extension. If $L_v$ is an unramified quadratic extension or is isomorphic to $L_w \oplus L_{w'}$ then $r_v$ is class 1 (a $K_v$-fixed vector being the characteristic function of the maximal order of $L_v$). Thus we get a representation $r$ of $\SL(2,\A(F))$ in $S_0(L_A) = \otimes S(L_v)$, a subspace of the Schwartz-Bruhat space $S(L_A)$, by setting

\begin{equation}\label{eq:7.35}
r(s)\left(\prod_v \Phi_v\right) = \prod_v r_v(L_v)(s_v)\Phi_v.\tag{7.35}
\end{equation}

The representation $r(s)$ is called the global Weil representation (attached to $L$ (and $\tau$)). It can be shown to extend to $S(L_A)$ in such a way that the mapping $(s, \Phi) \to r(s)\Phi$ of $\SL(2,\A) \times S(L_A)$ into $S(L_A)$ is continuous. Not surprisingly it also extends to a unitary representation in $\bL^2(L_A)$ such that

\begin{equation}\label{eq:7.36}
r\left(\begin{pmatrix} 0 & 1 \\ -1 & 0 \end{pmatrix}\right)\Phi(x) = \gamma\hat{\Phi}(x^{\sigma})\tag{7.36}
\end{equation}

and

\begin{equation}\label{eq:7.37}
r\left(\begin{pmatrix} 1 & u \\ 0 & 1 \end{pmatrix}\right)\Phi(x) = \tau(bN(u))\Phi(x)\tag{7.37}
\end{equation}

with $\gamma = \prod \gamma_v$

Now (at last) we can explain what this representation has to do with theta-series and the explicit construction of cusp forms. In \cite{Weil3} it is shown that the series

\begin{equation}\label{eq:7.38}
\Theta(\Phi) = \sum_{\xi \in L} \Phi(\xi)\tag{7.38}
\end{equation}

converges uniformly on any compact subset of $S(L_A)$. Therefore the function $\Theta(r(s)\Phi)$, for each $\Phi$ in $S(L_A)$, is a continuous function on $\SL(2,\A)$ which is left $\SL(2,F)$ invariant by virtue of \eqref{eq:7.36} and \eqref{eq:7.37}. The resulting space of functions

\begin{equation}\label{eq:7.39}
\Theta(r(s)\Phi) = \sum_{\xi \in L} r(s)\Phi(\xi)\tag{7.39}
\end{equation}

comprises a space of generalized theta functions.

To obtain cusp forms on $\GL(2)$ (actually $\SL(2)$) one decomposes this space of theta functions according to the characters $\lambda$ of the relative norm one group $H^1$ of $L_A$ (which are trivial on $H_F^1$). This decomposition runs parallel to the decomposition of the local Weil representations $r_v(L_v)$ according to the local characters $\lambda_w$.

More precisely, consider the functions

\begin{equation}\label{eq:7.40}
\Theta_{\Phi}(\lambda)(s) = \int_{H^1 \setminus H_A^1} \lambda(n)\left(\sum_{\xi \in L} r(s)\Phi(\xi n)\right) \dd n\tag{7.40}
\end{equation}

obtained by projecting the theta-functions $\Sigma r(s)\Phi(\xi)$ onto the space of theta-functions which transform under $L_A^1$ according to $\lambda$. These functions are cusps forms on $\SL(2)$ and their space realizes the (essentially) irreducible unitary representation $\pi(\lambda, V) = \otimes r_v(\lambda_v, L_v)$, the components $r_v(\lambda_v, L_v)$ of $r_v$ being defined in the obvious way.

In Section~\ref{ch:10} we shall describe a similar construction of cusp forms for $\GL(2)$ associated to the Weil representation of a global quaternion algebra $D$. In this case the theta functions that arise correspond to quaternary quadratic forms. Since the relation between these theta functions and certain classical theta series will be discussed in more detail then, we shall refrain now from making any further remarks about the theta-functions described by \eqref{eq:7.40} except to say that they generalize the theta-series with grossencharacter considered in \cite{Maass}.

\section{An Explicit Example}\label{sec:7c}

It has already been remarked several times that very little is known concerning which continuous series representations of $\SL(2,\R)$ occur discretely in $\bL^2(\Gamma_0(N)\setminus\SL(2,\R))$. We do know that the sequence of real numbers $\{s_j\}$ parameterizing such representations indexes non-zero spaces of real analytic cusp forms. However, the arithmetic significance of this sequence is not known.

The purpose of this subsection is to apply the results of Jacquet-Langlands theory to obtain explicit examples of such representations for various congruence subgroups. Although we shall supply details for but one example it will be clear from the context how to construct others. Unfortunately this construction tells us little about the nature of the sequence $\{s_j\}$ for a fixed given level $N$.

Let $L$ denote a quadratic extension of $\Q$. Our goal is to construct a grossencharacter $\lambda$ of $L$ such that $\pi(\lambda)$ is a cusp form for $\GL(2,\A(\Q))$ with the property that its infinite component is equivalent to some continuous series representation. This means that our quadratic extension should be real and our grossencharacter should not come from a grossencharacter of $\Q$.

To simplify matters we fix $L$ equal to $\Q(\sqrt{2})$. With this choice of discriminant only the prime $p = 2$ ramifies in $L$. We also want our grossen-character to be as unramified as possible. Therefore we set

\begin{equation}\label{eq:7.41}
\lambda((a_p)) = \prod_{j=1}^{2} |a_{\infty_j}|^{it_j} \prod_{p\ \mathrm{finite}} |a_p|_p^{is_p}\tag{7.41}
\end{equation}

where $\infty_j$ denotes the place at infinity corresponding to the embedding

\[
a + b\sqrt{2} \to a + (-1)^{j+1}b\sqrt{2}.
\]

(The constraint to keep in mind is that $\lambda$ must not arise from a grossen-character of $\Q$.)

Now the function defined by \eqref{eq:7.41} is clearly a character of the idele group $A^{\times}$ of $L$. It will be a grossencharacter of $L$ precisely when

\begin{equation}\label{eq:7.42}
\lambda(a) = 1\tag{7.42}
\end{equation}

for all principal id\`eles $\alpha$. This last condition imposes certain conditions on $t_1$ and $t_2$ that we shall now describe.

Let $I_S$ denote the set of $S$ units of $L$ with $S$ the set of infinite primes. Then $I_S$ is the group of units in the ring of integers of $L$ and is generated (in this case) by $-1$ and the fundamental unit $\sqrt{2} + 1$. Since \eqref{eq:7.42} must hold in particular for these $S$-units we must have

\begin{equation}\label{eq:7.43}
\lambda(-1) = 1\tag{7.43}
\end{equation}

and

\begin{equation}\label{eq:7.44}
\lambda(\sqrt{2} + 1) = 1.\tag{7.44}
\end{equation}

This first equation puts no restriction on $t_1$ and $t_2$, but the second implies that

\[
|\sqrt{2} + 1|^{it_1}|-\sqrt{2} + 1|^{it_2} = 1
\]

or, taking logarithms, that

\begin{equation}\label{eq:7.45}
it_1\log|\sqrt{2} + 1| + it_2\log|\sqrt{2} - 1| = 2k\pi i\tag{7.45}
\end{equation}

with $k \in \Z$. We note that infinitely many pairs $\{t_1, t_2\}$ may satisfy \eqref{eq:7.45}. However, since $(\sqrt{2} + 1) = (\sqrt{2} - 1)^{-1}$, each pair $t_1, t_2$ must be chosen such that $i(t_2 - t_1)\log(\sqrt{2} - 1) = 2k\pi i$, i.e.\ such that

\begin{equation}\label{eq:7.46}
s = t_2 - t_1 = \frac{2k\pi}{\log(\sqrt{2} - 1)}, \quad k \in \Z.\tag{7.46}
\end{equation}

Henceforth we fix $k = 1$ in \eqref{eq:7.46} and we assume $t_1 + t_2 = 0$. Having chosen fixed $t_1, t_2$ satisfying \eqref{eq:7.46} there will be at least one choice of $\{s_p\}$ such that the formula \eqref{eq:7.41} defines a grossencharacter of $L$. In fact since the class number of $\Q(\sqrt{2})$ is 1 there will be exactly one choice of such $\{s_p\}$. (Cf.\ the discussion in [Tate], p. 343.) Note that the resulting grossencharacter $\lambda$ can not possibly be of the form $\mu \circ N$ with $\mu$ a grossencharacter of $\Q$ since our choice of $k$ in \eqref{eq:7.46} implies $t_1 \neq t_2$.

From this discussion we conclude that the representation

\begin{equation}\label{eq:7.47}
\pi(\lambda) = \pi_{it_1,it_2} \otimes \left(\bigotimes_{p < \infty} \pi_p\right)\tag{7.47}
\end{equation}

occurs in $\bR_0^\psi$ if $\psi$ denotes the central character of $\pi(\lambda)$ and $\pi_p$ is the representation of $\GL(2,\Q_p)$ associated to $\lambda_p$ (or the pair $\lambda_p, \lambda_{p'}$ if both $p$ and $p'$ divide $p$). Note that $\pi_{it_1,it_2}$ is trivial on the center of $\GL(2,\R)$ since $t_1 + t_2 = 0$. Thus finally we can prove:

\begin{theorem}\label{thm:7.12}

The continuous series representation $\pi_{is}$, with

\[
s = \frac{2\pi}{\log(\sqrt{2} - 1)},
\]

occurs discretely in $\bL^2(\Gamma_0(2)\setminus\SL(2,\R))$.

\end{theorem}

\begin{proof}

We claim that the conductor of $\pi(\lambda)$ is 2. Indeed since only $p = 2$ ramifies in $L$, and $\lambda_p$ is unramified for every $p$, $\pi_p$ is class 1 for all odd $p$ (by the proof of Proposition~\ref{prop:7.8}). Thus the conductor of $\pi(\lambda)$ is the conductor of $\pi_2$.

To compute the conductor of $\pi_2$ let $\delta_2$ denote the character $|\cdot|_p^{is_p}$ where $p = \sqrt{2}$ lies above 2. Let $\omega_2$ denote the character (of order 2) of $\Q_2^{\times}$ attached to the quadratic extension $\Q_2(\sqrt{2})$. Then $\pi_2 = \pi(\delta_2, \delta_2\omega_2)$.

(This follows from the fact that $|a_p|_p^{is_p} = |N(a_p)|_{\Q_2}^{is_p}$ if $p = \sqrt{2}$.) So since $\omega_2$ has conductor 2 (or $2\mathcal{O}_2$) our claim is an immediate consequence of \eqref{eq:4.20}.

To prove the theorem let $H(t_1, t_2)$ denote the subspace of the representation space of $\pi(\lambda)$ in $\bL_0^2(\bG_\Q \setminus \bG_\A, \psi)$ consisting of functions which transform according to $\psi_2$ under the action of $K_0^2$. By our claim above $H(t_1, t_2)$ is equivalent (as a $\GL(2,\R)$-module) to $\pi_{it_1,it_2}$. But the restriction of $\pi_{it_1,it_2}$ to $\SL(2,\R)$ is irreducible and equivalent to $\pi_{is}$ where $s = t_1 - t_2$. (Cf.\ the ``Notes and References'' in Section~\ref{ch:4}.) Therefore, since

\[
Z_\infty^+ \bG_\Q \setminus \bG_\A / K_0^N \cong \Gamma_0(2)\setminus\SL(2,\R),
\]

$H(t_1,t_2)$ is equivalent as an $\SL(2,\R)$-module to an irreducible subspace of $\bL^2(\Gamma_0(2)\setminus\SL(2,\R))$ belonging to $\pi_{is}$.

\end{proof}

\section{Connections with Class Field Theory}\label{sec:7d}

Here we shall be vague as well as brief. Our remarks are directed toward the reader who would like to have some idea of the relationship between Jacquet-Langlands' theory and (possible future developments in) class field theory.

Let $F$ denote a global field and $L$ a Galois extension of $F$.

The central concern of class field theory is then the arithmetic study of such extensions. In particular one should obtain a description of how the $F$-primes split in $L$ and a ``reciprocity law'' describing the Galois group of this extension in terms of $F$.

A Galois extension is said to be abelian if its Galois group is abelian. In the case of abelian extensions the fundamental reciprocity law may be described as follows. Let $\mathfrak{G} = G(L/F)$ denote the Galois group of the abelian extension $L/F$ and let $\chi$ denote a character of $\mathfrak{G}$. Then one can canonically associate to $\chi$ a character of the group $I_F(d_F)$ (which is again denoted by $\chi$). (Here $I_F(d_F)$ is the group of $F$-ideals relatively prime to the discriminant $d_F$ of $F$; if $v$ is an $L$-unramified finite prime of $F$ and $w$ is a fixed prime of $L$ dividing $v$ one can set $\chi(v)$ equal to $\chi$ evaluated at the Frobenius automorphism at $v$ and extend $\chi$ to $I_F(d_F)$ by multiplicativity.) The ``abelian'' L-function corresponding to $\chi$ is

\begin{equation}\label{eq:7.48}
\begin{aligned}
L(s,\chi; L/F) &= \prod_{v\ \mathrm{unramified}} (1 - \chi(v)(Nv)^{-s})^{-1}\\
&= \sum_{\mathfrak{A}} \chi(\mathfrak{A})N(\mathfrak{A})^{-s}
\end{aligned}\tag{7.48}
\end{equation}

the summation extending over all ideals in $I_F(d_F)$.

On the other hand suppose $\psi$ is a grossencharacter of $F$, i.e.\ a unitary character of $\GL(1,\A(F)) = \A(F)^{\times}$ trivial on $\GL(1,F) = F^{\times}$. Suppose also that $\psi = \prod \psi_v$ and that $\psi_v$ is unramified for each $v$ not dividing $d_F$. This $\psi$ defines a ``Hecke character'' of $I_F(d_F)$ through the formula

\begin{equation}\label{eq:7.49}
\chi(v) = \psi_v(\tilde{\omega}_v)\tag{7.49}
\end{equation}

and the fundamental reciprocity law of abelian class field theory may be formulated as the assertion that every character $\chi$ of $\mathfrak{G}$ (viewed as a character of $I_F(d_F)$) arises in this way. That is, every character $\chi$ of $\mathfrak{G}$ is associated to some grossencharacter $\psi$ of $F$.

Now each grossencharacter $\psi$ of $F$ has associated to it a Hecke L-function $L(s,\psi)$. Therefore the reciprocity law above may alternately be written in the form

\begin{equation}\label{eq:7.50}
L(s,\chi; L/K) = L(s,\psi)\tag{7.50}
\end{equation}

where $\psi$ is some appropriate grossencharacter of $F$. Note that the left side of \eqref{eq:7.50} is an Euler product whose $v$-th factor depends on how $v$ splits in $L$. But the right side of \eqref{eq:7.50} is an Euler product whose $v$-factor depends only on $\psi_v$. Thus the arithmetic significance of an equality such as \eqref{eq:7.50} is evident.

Summing up, the fundamental result of abelian class field theory asserts the following. Suppose $\chi$ is any character of the Galois group of an extension $L/F$. Then there is a canonical correspondence which associates to $\chi$ an automorphic form $\psi(\chi)$ on $\GL(1)$ (over $F$) such that

\begin{equation}\label{eq:7.51}
L(s,\chi; L/K) = L(s,\psi(\chi)).\tag{7.51}
\end{equation}

This result may also be expressed without explicit reference to the extension field $L$ by considering in place of $L$ the maximal abelian extension $F_{ab}$ of $F$. If $G_{ab}$ denotes the Galois group of $F_{ab}/F$ one may assert that the dual of $G_{ab}$ is isomorphic to the torsion subgroup of the dual of the idele class group of $F$.

Now if we want to consider not necessarily abelian extensions of $F$ it is natural to introduce the Galois group $\mathfrak{G}$ of the algebraic closure of $F$ over $F$. A fundamental problem of non-abelian class field theory then is to describe the set $\mathfrak{G}^*(n)$ consisting of equivalence classes of irreducible unitary representations of $\mathfrak{G}$ of dimension $n$. Note that $\mathfrak{G}^*(n)$ is simply the dual group of $\mathfrak{G}_{ab}$ in case $n = 1$. (This is not to say that $\mathfrak{G} = \mathfrak{G}_{ab}$; rather every character of $\mathfrak{G}$ is trivial on the commutator subgroup of $\mathfrak{G}$ and hence defines a character of $\mathfrak{G}_{ab}$.) In case $n = 1$ the solution to the problem is therefore given by abelian class field theory.

If $n = 2$, let $X$ denote the collection of automorphic cusp forms on $\GL(2)$ over $F$, i.e.\ irreducible unitary representations $\pi$ of $\GL(2,\A(F))$ which occur in $\bR_0^\psi$ for some appropriate $\psi$. Then what one would hope to be able to do is define an injection from $\mathfrak{G}^*(2)$ into $X$ and describe its image completely. (This is analogous to specifying the torsion subgroup of the dual of $F^{\times}\setminus\A(F)^{\times}$ in case $n = 1$.)

One general result along these lines is due to Jacquet-Langlands. Before describing it let us point out that Theorem~\ref{thm:7.11} of Subsection B is more or less equivalent to the following. Suppose $\pi'$ is a monomial element of $\mathfrak{G}^*(2)$, i.e.\ $\pi'$ is a representation induced from a suitable character $\lambda$ of a subgroup of $\mathfrak{G}$ of index 2. Then one may associate to $\pi'$ an automorphic cusp form $\pi(\pi')$ on $\GL(2)$. This is so because the Artin L-function $L(\pi', s)$ is (in this case) entire and coincides with the L-function $L(s, \pi(\pi'))$. (Note that $\lambda$ may be identified with a character of the idele group of a quadratic extension of $F$, whence Theorem~\ref{thm:7.11}.

The result of Jacquet-Langlands alluded to above represents a vast generalization of Theorem~\ref{thm:7.11} and to state it properly requires additional terminology and reformulation.

If $F$ is a local field let $C_F$ denote the multiplicative group of $F$, and if $F$ is a global field, let $C_F$ denote the idele class group of $F$. If $L$ is a finite Galois extension of $F$ then the Weil group $W(L/F)$ is a certain extension of $\mathfrak{G} = G(L/F)$ by $C_L$. In particular $W(L/F) = C_F$ if $L = F$ and $W(L/F)/C_L = G(L/K)$ in general. If $F$ is global, $v$ is a place of $F$, and $w$ in $L$ divides $v$, there is a uniquely determined homomorphism $\alpha_v$ of $W(L_w/F_v)$ into $W(L/F)$ (uniquely determined up to inner automorphism by an element of $C_L$).

If $F$ is a local field then to every representation $\sigma$ of $W(L/F)$ one can associate a local ``L-function'' $L(s,\sigma)$ and local factor $\varepsilon(s,\sigma)$ both of which depend only on the class of $\sigma$. If $F$ is global and $\sigma$ is a semi-simple (complex) representation of $W(L/F)$ then $\sigma_v = \sigma \circ \alpha_v$ is a representation of $W(L_w/F_v)$ whose class is determined by that of $\sigma$. Therefore we set

\begin{equation}\label{eq:7.52}
L(s,\sigma) = \prod_v L(s,\sigma_v).\tag{7.52}
\end{equation}

The function defined by this product (in a certain half-plane) extends to a meromorphic function in all of $\C$ and satisfies the functional equation

\begin{equation}\label{eq:7.53}
L(s,\sigma) = \varepsilon(s,\sigma)L(1 - s,\tilde{\sigma})\tag{7.53}
\end{equation}

where $\tilde{\sigma}$ is contragredient to $\sigma$ and $\varepsilon(s,\sigma) = \prod_v \varepsilon(s,\sigma_v)$.

Now suppose $\sigma$ is a two dimensional representation of $W(L/F)$. From Theorem~\ref{thm:6.14} it follows that there is at most one irreducible admissible representation $\pi$ of $\mathcal{H}(G = \GL(2,F))$ such that $\pi\left(\begin{pmatrix} a & 0 \\ 0 & a \end{pmatrix}\right) = \det(\sigma(a))I$ and such that for all characters $\chi$ of $F^{\times} = C_F$,

\[
L(s,\chi \otimes \pi) = L(s,\chi \otimes \sigma)
\]

\[
L(s,\chi^{-1} \otimes \tilde{\pi}) = L(s,\chi^{-1} \otimes \tilde{\sigma})
\]

and

\[
\varepsilon(s,\chi \otimes \pi) = \varepsilon(s,\chi \otimes \sigma).
\]

(Here $\chi$ lifts to a one-dimensional representation of $W(L/F)$ through the natural homomorphism from $W(L/F)$ onto $W(F/F) = C_F$.) We note that $\pi = \pi(\lambda)$ if $L$ is a separable quadratic extension of $F$ and $\sigma$ is the representation of $W(L/F)$ induced from the character $\lambda$ of the subgroup $C_L = W(L/L)$. If $\sigma$ is the direct sum of one dimensional representations of $W(L/F)$ lifted (as above) from the characters $\mu_1$ and $\mu_2$ of $F^{\times}$ then $\pi(\sigma) = \pi(\mu_1, \mu_2)$. In general, $\pi$, if it exists at all, will be denoted by $\pi(\sigma)$.

\begin{theorem}\label{thm:7.13}

Suppose $F$ is a global field and $\sigma$ is a two-dimensional semi-simple representation of the Weil group of the separable quadratic extension $L$ of $F$. Let $\psi = \det(\sigma)$. Suppose also that for every character $\chi$ of $C_F$ both $L(s,\chi \otimes \sigma)$ and $L(s,\chi^{-1} \otimes \tilde{\sigma})$ are entire functions of $s$ bounded in vertical strips of finite width. Then $\pi(\sigma_v)$ exists for every place $v$ and $\pi(\sigma) = \otimes \pi(\sigma_v)$ is a representation of $\GL(2,\A(F))$ which occurs in $\bR_0^\psi$.

\end{theorem}

Observe that by the above remarks this result includes Theorem~\ref{thm:7.11} as a special case. (Take $\sigma$ equal to the representation of the global Weil group induced from the character $\lambda$ of $W(L/L)$. Then $\sigma_v = \Ind(\lambda_w)$ if $v$ does not split in $L$ but lies under $w$, and is the direct sum of two characters if $v$ does split in $L$.) In this case $L(s,\sigma)$ is known to be entire by the abelian theory. On the other hand, for function fields $F$ and arbitrary representations $\sigma$, the L-functions $L(s,\chi \otimes \sigma)$ are also known to be entire. Therefore Theorem~\ref{thm:7.12} is a ``real'' Theorem in this case as well.

For arbitrary fields it should be remarked that the converse to Theorem~\ref{thm:7.12} follows immediately from the ``only if'' part of Theorem~\ref{thm:6.18} since

\begin{equation}\label{eq:7.54}
L(s,\sigma) = L(s,\pi(\sigma)).\tag{7.54}
\end{equation}

That is, if $\pi(\sigma_v)$ can be shown to exist for every place $v$, and if $\bigotimes_v \pi(\sigma_v)$ can be shown to be a cusp form, then $L(s,\sigma)$ will be entire! Although this hardly constitutes a feasible attack on the conjecture that $L(s,\sigma)$ is always entire it does demonstrate (once again) the interest of Theorem~\ref{thm:6.18}.

Finally we observe that the \eqref{eq:7.54} is a natural generalization of the abelian reciprocity law \eqref{eq:7.51}.

\section*{Notes and References}

\addcontentsline{toc}{subsection}{Notes and References}

The application of Weil's representation to the construction of irreducible representations of other semi-simple groups is the subject matter of \cite{Gelbart4} and \cite{GrossKunze}. The example of Subsection C is taken from \cite{Gelbart2}.

Several important conjectures and methods related to Langlands' general ``philosophy of L-functions'' are to be found in \cite{Langlands2}. These conjectures are motivated by the results of \cite{Langlands3} and \cite{JacquetLanglands} represents an attempt to verify them for $\GL(2)$.

If one forgets about L-functions (for the moment) it still seems to be of interest to formulate the following:

\textbf{Problem A.} Given two groups $G'$ and $G$, what natural relations are there (if any) between automorphic forms on $G'$ and $G$?

Here we are assuming that $G$ and $G'$ are such that the notion of automorphic form on $G$ or $G'$ is meaningful. If $G$ is an algebraic group defined over a global field $F$ then an automorphic form must be an irreducible representation of the adele group of $G$ which occurs in the obvious $\bL^2$ space. An important example is when $G = \GL(2)$ (over $F$) and $G'$ is the multiplicative group of quadratic extension of $F$. In this case a solution to Problem A is provided by Theorem~\ref{thm:7.11} (and this solution is natural because the appropriate L-functions coincide).

In general one would hope that a solution to Problem A would make it possible to construct automorphic forms on certain groups using automorphic forms on simpler groups and to obtain arithmetic information from the explicit relations between these forms. Such is the case when $G = \GL(2)$ and $G'$ is the multiplicative group of a quaternion algebra (see Section~\ref{ch:10}). Other examples include those which appear in \cite{Jacquet} (where $G' = \GL(2)$ over $F$ and $G = \GL(2)$ over a quadratic extension of $F$) \cite{Gelbart3} (where $G'$ is (an extension of) Weil's metaplectic group and $G = \GL(2)$), and unpublished work of Shintani.

Our discussion in Subsection D is adapted from \cite{Shalika}, \cite{Langlands2} and \cite{JacquetLanglands}. The construction (and the motivation for the construction) of the groups $W(L|F)$ is explained in \cite{Weil5}. Weil also discusses the L-functions $L(s,\sigma)$ systematically developed in \cite{Langlands3}. These L-functions generalize the ``non-abelian'' L-series introduced in \cite{Artin} as well as the ``abelian'' L-series of Hecke.
